% !TEX root = ../frode-bibelen.tex
% ----------------------------------------------------------------------------
%  Frode-bibelen -- 5 Tekster
%  Inkluderes fra frode-bibelen.tex (kompiler alltid hovedfilen).
% ----------------------------------------------------------------------------

% ========================================================= 5 TEKSTER =======
\kapittel{Tekster (\texorpdfstring{\cc{char}}{char}-arrayer og \texorpdfstring{\cc{char*}}{char*})}

En tekst er en \cc{char}-array som slutter med \cc{'\bs 0'}. \cc{"Ola"} tar \textbf{4} plasser: \cc{'O' 'l' 'a' '\bs 0'}.

\begin{multicols}{2}
\textbf{Fast buffer} (lokal, kortvarig)
\begin{kode}
char buffer[STRLEN];
lesText("\tNavn", buffer, STRLEN);
\end{kode}
\columnbreak
\textbf{Dynamisk} (i structen, varig)
\begin{kode}
char* navn;                  //  Bare en PEKER!
navn = lagOgLesText("\tNavn");
...
free(navn);
\end{kode}
\end{multicols}

\subsection{\texorpdfstring{\cc{string.h}}{string.h} og \texorpdfstring{\cc{ctype.h}}{ctype.h}}
\begin{center}
\small\renewcommand{\arraystretch}{1.08}
\begin{tabular}{@{}L{4.6cm}L{11.0cm}@{}}
\toprule
\textbf{Funksjon} & \textbf{Gjør / returnerer} \\
\midrule
\cc{strlen(s)} & Antall tegn \textbf{uten} \cc{'\bs 0'}. \cc{strlen("Ola")} = 3. \\
\cc{strcpy(til, fra)} & Kopierer \cc{fra} inn i \cc{til} (overskriver). \cc{til} må ha plass! \\
\cc{strncpy(til, fra, n)} & Kopierer maks \cc{n} tegn (legger \emph{ikke} alltid på \cc{'\bs 0'}). \\
\cc{strcat(til, fra)} & Skjøter \cc{fra} på slutten av \cc{til}. \\
\cc{strcmp(a, b)} & \textbf{0 hvis like}, $<0$ hvis \cc{a} kommer før \cc{b} alfabetisk, $>0$ ellers. \\
\cc{strncmp(a, b, n)} & Som \cc{strcmp}, men ser bare på de \cc{n} første tegnene. \\
\cc{strchr(s, c)} & Peker til \textbf{første} forekomst av tegnet \cc{c} i \cc{s}, ellers \cc{NULL}. \\
\cc{strrchr(s, c)} & Peker til \textbf{siste} forekomst av \cc{c}, ellers \cc{NULL}. \\
\cc{strstr(s, del)} & Peker til første forekomst av teksten \cc{del} i \cc{s}, ellers \cc{NULL}. \\
\midrule
\cc{toupper(c)}, \cc{tolower(c)} & Stor/liten bokstav (andre tegn uendret). \\
\cc{isupper}, \cc{islower}, \cc{isalpha}, \cc{isdigit}, \cc{isspace} & Sann (ikke 0) hvis tegnet er stor/liten bokstav, bokstav, siffer, blank. \\
\bottomrule
\end{tabular}
\end{center}

\subsection{Mønstre du trenger}
\begin{kode}
//  1) Er to tekster LIKE?   (ALDRI  a == b  -  det sammenligner adresser!)
if (strcmp(a, b) == 0)  ...          //  eller:  if (!strcmp(a, b))
//  2) Inneholder navnet en deltekst?   (søk på del av navn)
if (strstr(gKunder[i]->navn, sok))  ...      //  != NULL  =>  funn
//  3) Kopiere fra buffer til NY memory med eksakt riktig størrelse:
kunde->navn = (char*) malloc((strlen(buffer) + 1) * sizeof(char));   //  +1 for '\0'
strcpy(kunde->navn, buffer);
//  4) Gå gjennom tegn for tegn:
for (int i = 0;  i < strlen(t);  i++)   if (t[i] == 'a')  antall++;
char* p = t;   while (*p != '\0')  { if (*p == 'a') antall++;  p++; }
//  5) Bygge opp en ny tekst tegn for tegn - HUSK avsluttende '\0':
int n = 0;
for (int i = 0;  t[i];  i++)   if (isalpha(t[i]))  ny[n++] = toupper(t[i]);
ny[n] = '\0';
//  6) Alle forekomster av et tegn med strchr:
char* p = strchr(t, 'e');
while (p != NULL)  {  antall++;   p = strchr(p+1, 'e');  }
//  7) Tegn <-> tall:
int siffer = tegn - '0';       //  '7' -> 7
char c     = '0' + 5;          //  5   -> '5'
int nr     = toupper(c) - 'A'; //  'C' -> 2  (plass i alfabetet)
\end{kode}

\begin{felle}[{Tekst-feller}]
\begin{itemize}
  \item \cc{malloc(strlen(t))} er \textbf{én for lite} -- husk \cc{+1} til \cc{'\bs 0'}.
  \item \cc{strcpy} til en \cc{char*} som ikke peker på noe (ikke \cc{malloc}-et) krasjer.
  \item \cc{fgets} \textbf{beholder} \cc{'\bs n'}: fjern med \lstinline|t[strlen(t)-1] = '\0';|
  \item \cc{sizeof(peker)} er pekerens størrelse (8), ikke tekstens lengde. Bruk \cc{strlen}.
  \item \cc{printf("\%s", t+3)} skriver teksten \textbf{fra indeks 3} og ut -- nyttig i oppgave 1.
\end{itemize}
\end{felle}
